\begin{ExerciseBlock}
\ExerciseHeading{07}{Strahlungsabsorption in Aluminium}{geometrical_and_wave_optics_aluminium_absorption}
Röntgenstrahlung durchdringt eine Aluminiumprobe. Ihre Intensität folgt dem Beerschen Absorptionsgesetz. Vor der Probe beträgt sie $I_0=\qty{100}{\watt\per\metre\squared}$, nach einer Schicht der Dicke $D=\qty{0.5}{\centi\metre}$ noch $I=\qty{0.34}{\watt\per\metre\squared}$.
\begin{enumerate}
\item Berechnen Sie den Absorptionskoeffizienten $\mu_\mathrm{Al}$ in \unit{\per\milli\metre}.
\item Bestimmen Sie die Halbwertsdicke $d_\mathrm H$.
\end{enumerate}
\end{ExerciseBlock}
\UseExerciseSolution{geometrical_and_wave_optics_aluminium_absorption}
